%%==============================================================================
%% Section I: Introduction
%%==============================================================================
\section{Introduction}
\label{sec:introduction}

\begin{figure}[t]
  \centering
  \begin{tikzpicture}[font=\tiny,
    wnode/.style={circle,inner sep=0pt,minimum size=2.4mm,fill=WorkerBlue!30,draw=WorkerBlue},
    hnode/.style={rectangle,rounded corners=2pt,inner sep=1.5pt,minimum height=5mm,
                  minimum width=13mm,fill=HeadTeal!15,draw=HeadTeal},
    snode/.style={rectangle,rounded corners=3pt,inner sep=2pt,minimum height=7mm,
                  minimum width=34mm,fill=ServerRed!12,draw=ServerRed},
    flow/.style={-{Latex[length=1.5mm]},gray!75,line width=0.35mm}]
    % --- Cloud ---
    \node[snode] (S) {\scriptsize Cloud server $S$};
    % --- Cluster heads ---
    \foreach \i/\x in {0/-20,1/0,2/20}{
      \node[hnode] (H\i) at ($(S.south)+( \x mm,-11mm)$) {\scriptsize CH $c_{\i}$};
    }
    % --- Workers ---
    \foreach \i in {0,1,2}{
      \foreach \k/\dx in {0/-4.5,1/0,2/4.5}{
        \node[wnode] (w\i\k) at ($(H\i.south)+( \dx mm,-6.5mm)$) {};
      }
    }
    % --- Links ---
    \foreach \i in {0,1,2}{
      \draw[flow] (H\i) -- (S);
      \foreach \k in {0,1,2}{\draw[flow] (w\i\k) -- (H\i);}
    }
    % --- Tier labels ---
    \node[HeadTeal,align=center,font=\tiny] at ($(H1)+(0,-1.5mm)$) {};
    \node[anchor=north,font=\tiny,text=black!75] at ($(S.south)+(0,-20mm)$)
      {\textbf{Cluster-head UAVs (strong)}};
    \node[anchor=north,font=\tiny,text=black!75] at ($(S.south)+(0,-31mm)$)
      {\textbf{Worker UAVs (MCU-class)}};
  \end{tikzpicture}
  \caption{Three-tier hierarchical FL over a UAV swarm. MCU-class worker UAVs are grouped into clusters, each governed by a capable cluster-head (CH) that performs the cryptographic workload and pre-aggregation; the cluster-heads connect to a cloud server that performs global aggregation and coordinates verification.}
  \label{fig:hierarchy}
\end{figure}

\IEEEPARstart{U}{nmanned} aerial vehicle (UAV) swarms are increasingly deployed for disaster response, environmental monitoring, reconnaissance, and intelligent surveillance, in which a fleet of small, battery-powered aircraft collaborates to sense and act over a target area. Training a shared perception or decision model across the fleet is essential, yet the raw imagery, telemetry, and location traces are both volume-heavy and privacy-sensitive, which makes centralizing them impractical. Federated learning (FL) offers a natural fit: each vehicle trains locally on its own data and exchanges only model updates, so that raw samples never leave the device~\cite{bonawitz2019towards}. Because a flat, fully connected topology does not scale to large fleets, practical deployments organize the aircraft into a three-tier hierarchy in which worker UAVs are grouped into clusters governed by a more capable cluster-head UAV, and the cluster-heads in turn connect to a cloud server. Hierarchical aggregation over this topology reduces both the number of long-range links and the energy cost incurred at the weakest nodes~\cite{liu2020hierfavg},~\cite{wang2021rflha}, and it has become the standard structure for UAV-based collaborative learning~\cite{ahmed2025uav},~\cite{nagarajan2025drone}.

Exchanging model updates, however, does not by itself preserve privacy. A broad line of gradient-inversion and reconstruction attacks demonstrates that, from a single update, an adversary can recover much of the underlying training data~\cite{zhu2019deep},~\cite{geiping2020inverting},~\cite{zhao2020idlg}, and recent studies show that these attacks remain effective against modern architectures~\cite{guo2025exploring}. In a swarm, the cloud and even peer cluster-heads occupy strong vantage points and are therefore realistic adversaries. Three families of countermeasures are commonly adopted. Differential privacy (DP)~\cite{mcmahan2017learning} injects calibrated noise but incurs a utility trade-off; secure aggregation built on secure multi-party computation~\cite{bonawitz2017practical} masks pairwise updates; and homomorphic encryption (HE)~\cite{gentry2009fully} allows the cloud to aggregate ciphertexts without learning any individual update. HE is particularly attractive for secure aggregation, and the Cheon--Kim--Kim--Song (CKKS) scheme~\cite{cheon2017ckks}, which is built on the ring learning-with-errors (RLWE) problem~\cite{lyubashevsky2010ideal}, supports approximate arithmetic and is well matched to real-valued neural-network updates.

Deploying HE on a swarm exposes a sharp resource asymmetry. The dominant cost of CKKS encryption is a number-theoretic transform over a large polynomial ring, whose memory and energy footprint are prohibitive for the MCU-class processors that power small worker UAVs, while the larger cluster-heads can perform it. Outsourcing encryption to a cluster-head is the natural remedy, but it immediately raises two new problems. First, the cluster-head and cloud that mediate the outsourced computation are themselves untrusted, so a plaintext update handed to them is exposed. Second, conventional HE relies on a globally shared key that would allow any one intermediary to decrypt every update; existing outsourcing protocols therefore introduce a key distribution center or a trusted authority, and some additionally rely on a trusted execution environment~\cite{li2025fedphe},~\cite{wu2025esafl},~\cite{yuan2026fedout},~\cite{cai2023secfed}. In fielded UAV missions, however, no such trusted infrastructure is available.

Removing the trusted anchor is not sufficient, because two structural vulnerabilities remain. On the one hand, intermediaries are not merely curious but potentially malicious: a cluster-head may encrypt a modified plaintext and thereby poison the global model, and the cloud may fabricate an aggregate, selectively discard updates, or request the decryption of a chosen individual. On the other hand, the threat model explicitly permits the cluster-heads and cloud to collude, which is fatal to a naive two-party outsourcing design in which one party sees the blinded plaintext and the other holds a single masking factor, since their two views reconstruct the update. Prior verifiable protocols either stop at decryption-time verification without an inclusion proof~\cite{bouamama2025vesafl}, or add a dedicated masking server that is not part of the fixed three-tier topology. Three challenges therefore emerge:

\begin{itemize}[leftmargin=*,topsep=2pt,itemsep=1pt]
  \item \textbf{C1: Key establishment without trust.} How can a group decryption key be established without a trusted authority when worker UAVs cannot participate in key generation and the cloud must not be able to decrypt alone?
  \item \textbf{C2: Entity-free collusion-resistant outsourcing.} How can encryption be outsourced to a cluster-head while hiding the update from the cluster-head and cloud even when they collude, using only the entities already present in the three-tier hierarchy?
  \item \textbf{C3: End-to-end verifiability and dropout tolerance.} How can encryption, aggregation, and decryption be made publicly verifiable against a malicious cloud and corrupted cluster-heads while tolerating the frequent worker and cluster-head dropouts of wireless swarms?
\end{itemize}

We address these questions with \textsc{VHD-FL}. To answer C1, the cluster-heads run a distributed key generation (DKG) to establish a group key, and the resulting secret is shared among the cluster-heads through a threshold verifiable secret-sharing scheme; worker UAVs take no part in key generation and the cloud holds no share, so the cloud alone cannot decrypt. To answer C2, each worker masks its update with a blinding factor split into three additive shares and delivers them to its own cluster-head, the cloud, and a neighboring cluster-head arranged in a ring. The neighboring cluster-head is an existing strong node, so no new entity is introduced, and the three shares form a 3-out-of-3 secret sharing that resists collusion between any two holders. To answer C3, \textsc{VHD-FL} enforces three-layer verifiability: a zero-knowledge proof binds each ciphertext to the blinded plaintext, a homomorphic hash detects a forged aggregate or selective omission, and a further proof validates every partial decryption; threshold reconstruction and zero ciphertexts tolerate dropouts.

The main contributions are summarized as follows:
\begin{itemize}[leftmargin=*,topsep=2pt,itemsep=1pt]
  \item We design a \textbf{trust-anchor-free cluster-head DKG} that integrates a threshold verifiable secret sharing into the aggregation key. Worker UAVs contribute nothing to key generation and the cloud holds no share, yielding collusion-resistant threshold decryption without a trusted authority.
  \item We propose \textbf{three-blind outsourced encryption} that reuses a neighboring cluster-head as the third mask holder. It realizes 3-out-of-3 additive secret sharing, resists collusion between any two mask holders without adding an entity, and reduces per-round worker work to a single modular subtraction plus one lightweight signature.
  \item We establish \textbf{three-layer verifiability}---encryption zero-knowledge proofs, a homomorphic hash over the aggregation, and decryption proofs---against a malicious cloud and corrupted cluster-heads. We prove the corresponding privacy, verifiability, and dropout-tolerance properties, and experiments on MNIST, CIFAR-10, and FEMNIST show that \textsc{VHD-FL} matches plaintext FedAvg accuracy while substantially reducing worker and communication overhead.
\end{itemize}

The remainder of this paper is organized as follows. Section~\ref{sec:related} reviews related work. Section~\ref{sec:preliminaries} introduces the system model, threat model, cryptographic building blocks, and notation. Section~\ref{sec:scheme} details the \textsc{VHD-FL} scheme, and Section~\ref{sec:algorithms} presents the algorithms. Section~\ref{sec:security} gives the security analysis, Section~\ref{sec:evaluation} reports the performance evaluation, and Section~\ref{sec:conclusion} concludes the paper and discusses future directions.
